\section{A normalized family and its formal Gaussian action}
\label{sec:linear}

The stationary argument begins with a family that preserves the types
of all ends.  Compact observations make its Gaussian linearization
invertible, while separate infinite-end estimates identify the allowed
cylindrical and conical behavior.  The resulting analytic family is the
base for successive polynomial opening jets and their formal Gaussian
action.  This section constructs that family and the finite real
comparison used later.  The additional estimates on finite complex
contours are developed with the lateral actions in \Cref{sec:arrays}.

Gaussian Fredholm theory and finite-dimensional analytic reduction
have precedents in \cite[Sections~3.4 and~4]{CS2021}.  The end estimates
here also accommodate cylindrical opening jets, whose coefficients need
not form a convergent series.

Throughout this section, $S$ is a smooth properly embedded shrinker with
finitely many smooth conical and round cylindrical ends.  The cylindrical
radius is $\sqrt2$.  The end structure gives bounded curvature, bounded
$X^\perp$, bounded geometry, and polynomial area growth.  Constants may
depend on $S$, on a fixed finite number of derivatives of its end charts,
and on the fixed angular differentiability order.  They do not depend on
the length of an exhausted end.  Dependence on other parameters will be
specified explicitly.

Put $\rho_G=\exp(-|X|^2/4)$, $d\gamma=\rho_G\,dA_S$, and
\[
 L=\Delta_S-\tfrac12X^\top\cdot\nabla_S+\tfrac12+|\mathrm{II}|^2.
\]
All complexifications below are bilinear complexifications of the real
geometric formulas.  Hermitian products are used only in estimates.
An observation means a real continuous linear functional
$P_i u=\langle u,\chi_i\rangle_{L^2_G}$ with
$\chi_i\in C_c^\infty(S)$.  The observations are defined in fixed graph
coordinates on their common compact support.

\subsection{The Gaussian domain and compact observations}

\begin{lemma}[Gaussian realization]\label{lin:gaussian}
Let $H^2_G$ be the completion of $C_c^\infty(S)$ for
\[
 \|u\|_{H^2_G}=\sum_{i=0}^2\|\nabla^iu\|_{L^2_G}.
\]
The self-adjoint realization of $L$
has domain $H^2_G$, compact resolvent, and
\begin{equation}\label{lin:gaussian-estimate}
 \|u\|_{H^2_G}+\||X|\nabla u\|_{L^2_G}
       +\||X|^2u\|_{L^2_G}
 \le C\bigl(\|Lu\|_{L^2_G}+\|u\|_{L^2_G}\bigr).
\end{equation}
In particular its kernel $K$ is finite dimensional and consists of smooth
functions.
\end{lemma}

\begin{proof}
The shrinker identities imply
$\operatorname{div}_S X^\top=2-|X^\perp|^2/2$.  Weighted integration of
$\operatorname{div}_G(u^2X^\top)$ therefore gives
\[
 \tfrac12\int_S |X|^2u^2\,d\gamma
 =2\int_Su^2\,d\gamma+2\int_SuX^\top\cdot\nabla u\,d\gamma.
\]
Cauchy's inequality proves
$\||X|u\|_G\le C(\|\nabla u\|_G+\|u\|_G)$.
The same calculation applied to the norm squared of $\nabla u$ gives
$\||X|\nabla u\|_G\le C(\|\nabla^2u\|_G+\|\nabla u\|_G)$.
Applying the first inequality to $(1+|X|^2)^{1/2}u$ then controls
$|X|^2u$.  These calculations initially use compact supports, and pass to
the completion.

For completeness the domain can also be read from an ordinary
Schr\"odinger operator.  The unitary map $Uu=\rho_G^{1/2}u$ satisfies
\begin{equation}\label{lin:half-density}
 -ULU^{-1}=-\Delta_S+\frac{|X|^2}{16}
       +\frac{|X^\perp|^2}{16}-1-|\mathrm{II}|^2.
\end{equation}
Indeed, for $\phi=|X|^2/4$ one has
$\nabla\phi=X^\top/2$ and $\Delta\phi=1-|X^\perp|^2/4$.
The last three terms in \eqref{lin:half-density} are bounded.
If $q=|X|^2/16$, direct integration by parts yields
\[
 \|(-\Delta+q)w\|_2^2
 =\|\Delta w\|_2^2+\|qw\|_2^2
       +2\int q|\nabla w|^2-\int(\Delta q)|w|^2.
\]
Here $\Delta q$ is bounded.  The Bochner identity and bounded Ricci
curvature control $\nabla^2w$ by $\Delta w$ and lower derivatives.
It follows that the closed form realization has domain
\[
 \{w:\nabla^2w,\ |X|\nabla w,\ |X|^2w,w\in L^2(S)\},
\]
with equivalent graph norm.  Cutoffs and local elliptic regularity show
that compactly supported smooth functions are a core.  Conjugating back,
and using the moment estimates just proved, identifies this domain with
the derivative-defined $H^2_G$ and proves
\eqref{lin:gaussian-estimate}.  Finally $q\to\infty$ on every end.
The $\|qw\|_2$ bound makes the $L^2$ tails uniformly small, and Rellich
compactness on compact subsets gives compactness of the domain inclusion.
\end{proof}

\begin{lemma}[Symmetric enlargement]\label{lin:observations}
There is a finite list of compact observations for which
\begin{equation}\label{lin:augmented}
 \mathcal A(u,\lambda)
    =\bigl(Lu+P^*\lambda,Pu\bigr):
 H^2_G\oplus\mathbb R^m\longrightarrow L^2_G\oplus\mathbb R^m
\end{equation}
is an isomorphism.  Moreover, let $u_1,\ldots,u_a\in H^2_G$ have
smooth compactly supported $Lu_j$.  The list can be enlarged while preserving
this isomorphism so that each $u_j$ occurs among the derivatives of the
normalized homogeneous solution family
$Lu+P^*\lambda=0$ obtained by varying $Pu$.
\end{lemma}

\begin{proof}
Choose $k=\dim K$ compact smooth functions $\chi_i$ whose pairings with a
basis of $K$ form a nonsingular matrix.  For this initial list, projection
of the first equation of \eqref{lin:augmented} onto $K$ determines
$\lambda$.  Inversion of $L$ on $K^\perp$ determines the component of
$u$ in $K^\perp$, and $Pu$ determines its component in $K$.

Let $G=(L|_{K^\perp})^{-1}$ and
$\mathcal V=C_c^\infty(S)\cap K^\perp$.  This space is dense in
$K^\perp$: approximate by compact functions and remove their finitely
many kernel pairings using the initial $\chi_i$.  The symmetric form
\[
 Q(f,g)=\langle Gf,g\rangle_G,\qquad f,g\in\mathcal V,
\]
is nondegenerate.  If $Q(f,g)=0$ for all $g\in\mathcal V$, density gives
$Gf=0$ and hence $f=0$.  The sources $q_j=Lu_j$ lie in $\mathcal V$,
by self-adjointness.

Every finite subspace $V\subset\mathcal V$ has a finite nondegenerate
enlargement for $Q$.  To see this, split off a nondegenerate part of $V$
and let $r_1,\ldots,r_s$ span its radical.  Nondegeneracy on
$\mathcal V$ supplies vectors $w_i$ with $Q(r_i,w_j)=\delta_{ij}$.
Subtract their projections to the already nondegenerate part.  On
$\operatorname{span}(r_i,w_i)$ the matrix is
$\left(\begin{smallmatrix}0&I\\I&C\end{smallmatrix}\right)$ and is
invertible.  Apply this construction to the span of the $q_j$ and add
the resulting compact functions as observation duals.

For the enlarged homogeneous block, kernel projection first forces the
original $k$ multipliers to vanish.  Its other multipliers vanish by the
nondegeneracy of $Q$ on the added space; the original observations then
kill the remaining kernel component of $u$.  The same decomposition
solves the inhomogeneous equations and proves bijectivity.  Finally
$Lu_j=q_j\in\operatorname{range}P^*$, so $u_j$ is the normalized
solution associated with the observation value $Pu_j$ and a suitable
multiplier vector.  No positivity of $Q$ is required.
\end{proof}

We use the sign convention in which a normalized stationary graph
satisfies
\begin{equation}\label{lin:normalized-equation}
 \Phi(u)+\sum_i\lambda_i\chi_i(u)=0,\qquad P(u)=q.
\end{equation}
Here the compact functions $\chi_i(u)$ include the smooth density and
normal-velocity factors converting the fixed observation into its
variational dual.  The equation is written in scalar normal-velocity
units, identifying $\Phi$ with its component in the chosen normal
direction.  Thus
$dF(u)[v]=\sum_i\lambda_i\,dP_i(u)[v]$ on a normalized solution.
At $(u,\lambda)=(0,0)$ the derivative of
\eqref{lin:normalized-equation} is \eqref{lin:augmented}.
Exterior scalar variables and drift normalizations below also identify
the range: whenever an equation is multiplied by a nonsingular scalar
factor, its source is multiplied by the same factor.  These are
compatible local descriptions of the geometric normalized block, not
an alteration of its compact variational normalization.

\subsection{Cylindrical strip estimates, traces, and separated rates}

On a round end write $z=e^t$, $\mu=e^{-t}$, and use the radius
$R(t,\theta)$ as graph variable.  The curvature numerator in these
coordinates is
\[
 g^{zz}R_{zz}+2g^{z\theta}
       (R_{z\theta}-R_zR_\theta/R)
       +g^{\theta\theta}(R_{\theta\theta}-R-2R_\theta^2/R),
 \qquad g=\operatorname{diag}(1,R^2)+dR\otimes dR.
\]
The support-function term is $(R-zR_z)/2$.  At $R=\sqrt2$, minus twice
the linearized radial graph velocity is the drift-normalized operator
\begin{equation}\label{lin:cylinder-operator}
 \mathcal L_{\mathrm{cyl}}v
 =v_t-2\mu^2(v_{tt}-v_t)-v_{\theta\theta}-2v.
\end{equation}
In particular the limiting slow rates are
$\sigma_j=2-j^2$, $j\in\mathbb Z$.  In circular data the equation with
the axial diffusion omitted is $R_t-R+2/R=0$, and the angular diffusion
coefficient is $2/R^2$.

Fix an integer $h\ge4$.  On a unit strip $I\times S^1$, with the
diffusion scale $E(t)$ comparable to $E_I\in(0,E_*]$, define
\begin{equation}\label{lin:strip-norm}
 \begin{split}
 N_{E_I,h}(v;I)^2={}&\|v\|_{L^2H^h}^2+\|v_t\|_{L^2H^h}^2
 +\|v_{\theta\theta}\|_{L^2H^h}^2\\
 &+E_I^2\|v_{tt}\|_{L^2H^h}^2
 +E_I\|v_{t\theta}\|_{L^2H^h}^2.
 \end{split}
\end{equation}
Second derivatives in this definition are bulk $L^2$ quantities.
The traces used below are traces of the function and its first
derivatives, and do not include a trace of $v_{tt}$.

\begin{lemma}[Strip and boundary estimates]\label{lin:strip}
Consider
\begin{equation}\label{lin:admissible-operator}
 \mathcal L v=v_t-a(t)v_{tt}-2b(t)v_{t\theta}
       -c(t)v_{\theta\theta}+d_1(t)v_t+d_2(t)v_\theta+d_0(t)v
\end{equation}
on an interval of real $t$ coordinates.  Suppose
\[
 \operatorname{Re}a\ge c_*E,\quad |a|\le C_*E,
 \quad\operatorname{Re}c\ge c_*,\quad |c|\le C_*,
 \quad |b|\le\varepsilon_*\sqrt E,
\]
the normalized derivatives of these coefficients are bounded,
$|d_1|\le C_*E$, and $d_0,d_2$ and their required derivatives are
bounded.  The leading coefficients are independent of $\theta$.
Sufficiently small operator-norm perturbations from the strip domain
to $L^2_tH^h_\theta$ are allowed.  The constants below depend only on
these bounds and a fixed strict accretivity margin.

There is a patch estimate
\begin{equation}\label{lin:patch}
 N_{E_I,h}(v;I)
 \le C\bigl(\|\mathcal Lv\|_{L^2(I^+;H^h)}
                   +\|v\|_{L^2(I^+;H^h)}\bigr),
\end{equation}
where $I^+$ is a fixed enlargement.  The corresponding boundary
estimate holds with homogeneous entrance Dirichlet or outgoing
derivative data.  Moreover
\begin{equation}\label{lin:first-traces}
 \sup_{t\in I}\bigl(\|v(t)\|_{H^h}
 +\|v_\theta(t)\|_{H^h}+\sqrt{E_I}\|v_t(t)\|_{H^h}\bigr)
 \le C N_{E_I,h}(v;I^+).
\end{equation}
For inhomogeneous boundary data one can use the spaces
\begin{align}
 \|\varphi\|_{\mathcal T^{D}_{E,h}}^2
 &=\sum_{j\in\mathbb Z}(1+j^2)^h(1+j^2)
                  (1+\sqrt E|j|)|\varphi_j|^2,
       \label{lin:dirichlet-trace}\\
 \|\psi\|_{\mathcal T^{N}_{E,h}}^2
 &=E\sum_{j\in\mathbb Z}(1+j^2)^h
                  (1+\sqrt E|j|)|\psi_j|^2.
       \label{lin:neumann-trace}
\end{align}
They are bounded by the bulk norm for $\varphi=v|_{\partial I}$ and
$\psi=v_t|_{\partial I}$.  In particular an outgoing derivative lift
has norm at most $C\sqrt E\|\psi\|_{H^{h+1}}$.
\end{lemma}

\begin{proof}
Freeze $a,c$ and first omit the small mixed coefficient.  For real
$\xi$ and integer $j$, accretivity and then the imaginary part give
\begin{equation}\label{lin:symbol}
 1+|i\xi+a\xi^2+cj^2|
 \ge c\bigl(1+|\xi|+E\xi^2+j^2+\sqrt E|j\xi|\bigr).
\end{equation}
Indeed the real part controls $E\xi^2+j^2$, their imaginary parts have
size at most a constant times that quantity, and
$2\sqrt E|j\xi|\le E\xi^2+j^2$.  The mixed term is absorbed by choosing
$\varepsilon_*$ small.  Fourier transformation in $t$ and Fourier
series in $\theta$ give the interior estimate.  The bounded normalized
derivatives allow freezing on intervals of fixed sufficiently small
length.  Commutators with a cutoff contain $a\chi'v_t$ and
$(\chi'-a\chi'')v$, so interpolation absorbs the derivative term and
leaves precisely the lower norm in \eqref{lin:patch}.

For the boundary estimate add a fixed positive zeroth-order shift
$K$.  The characteristic roots satisfy
\[
 a\lambda^2-\lambda-cj^2-K=0.
\]
There is one root in each open half-plane, and, with labels $+$ and $-$,
\begin{equation}\label{lin:roots}
 \begin{split}
 \operatorname{Re}\lambda_+\asymp|\lambda_+|
     &\asymp E^{-1}(1+\sqrt E|j|),\\
 -\operatorname{Re}\lambda_-\asymp|\lambda_-|
     &\asymp\frac{1+j^2}{1+\sqrt E|j|}.
 \end{split}
\end{equation}
To verify the half-plane assertion, an imaginary root would make the
real part of the equation equal to
$-\operatorname{Re}a\,\xi^2-\operatorname{Re}c\,j^2-K<0$.
Continue from positive real coefficients.  For $\sqrt E|j|$ small the
roots have the two asserted scales by the quadratic formula; for this
quantity large they are asymptotic to
$\pm j\sqrt{c/a}$, whose real parts have a uniform angle margin.
The intervening scaled compact set, together with the preceding
exclusion of imaginary roots, gives uniform constants.  These arguments
also cover the bounded transition between the two scales.

At an entrance the correction to a whole-line particular solution is
$\varphi_j e^{\lambda_-t}$.  At an exit, written as $t=0$ with $t<0$
inside, the derivative lift is
$\psi_j\lambda_+^{-1}e^{\lambda_+t}$.  Integrating their five squared
norms in \eqref{lin:strip-norm} gives, respectively,
\eqref{lin:dirichlet-trace} and \eqref{lin:neumann-trace}, up to uniform
constants.  This constructs the boundary parametrices.  Freezing and
absorption as above give the variable-coefficient boundary estimate;
removing $K$ only restores the lower norm in \eqref{lin:patch}.

The simpler traces follow by product integration, in particular from
\[
 \partial_t\|v_\theta\|_{H^h}^2
 =-2\operatorname{Re}\langle v_t,v_{\theta\theta}\rangle_{H^h},
 \qquad
 \partial_t\|v_t\|_{H^h}^2
 =2\operatorname{Re}\langle v_{tt},v_t\rangle_{H^h}.
\]
For the stronger Dirichlet trace apply the same argument in each mode
also to $\sqrt E|j|^3|v_j|^2$; its derivative is bounded by a constant
times $|j^2v_j|\,|\sqrt E jv_{t,j}|$.  For the Neumann trace set
$w=v_{t,j}$ and apply the one-dimensional trace inequality on length
$\ell=E(1+Ej^2)^{-1/2}$.  Multiplication by
$E\sqrt{1+Ej^2}$ bounds its right hand side by
\[
 C\bigl((1+Ej^2)\|v_{t,j}\|_2^2
                         +E^2\|v_{tt,j}\|_2^2\bigr).
\]
Summing proves the assertions.  Thus the additional angular regularity
in a convenient smooth-data lift is imposed on prescribed data, not
lost from the unknown.  A derivative trace supplied by another solution
already belongs to the natural space \eqref{lin:neumann-trace}.
\end{proof}

For a positive weight $\rho$ with bounded logarithmic derivatives, set
\begin{equation}\label{lin:weighted-strip}
 \|v\|_{X_{\rho,h}}
   =\sup_\nu\rho(t_\nu)^{-1}N_{E_\nu,h}(v;I_\nu),\qquad
 \|f\|_{Y_{\rho,h}}
   =\sup_\nu\rho(t_\nu)^{-1}\|f\|_{L^2(I_\nu;H^h)}.
\end{equation}
The intervals have length one and uniformly bounded overlap.
Equivalently one can put $\rho^{-1}v$ inside $N_{E_\nu,h}$.  If
$(\log\rho)'=d$ is a sufficiently large fixed positive number and
$E_*d$ is sufficiently small, the entrance Dirichlet/outgoing
derivative problem satisfies
\begin{equation}\label{lin:forward}
 \|v\|_{X_{\rho,h}}
 \le C\left(\|v(t_0)\|_{\mathcal T^D_{E(t_0),h}}
       +\|\mathcal Lv\|_{Y_{\rho,h}}
       +\rho(T)^{-1}\|v_t(T)\|_{\mathcal T^N_{E(T),h}}\right),
\end{equation}
where $\rho(t_0)=1$ and the constant is independent of $T$.
Here and below the logarithmic rate must exceed the needed propagation
rate by a fixed positive buffer.

To prove \eqref{lin:forward}, put $v=\rho w$.  For constant rate $d$ the
new mass is $d-ad^2+d_0+dd_1$.  Integration against $\bar w$ controls
$E|w_t|^2+|w_\theta|^2+d|w|^2$ if $d$ exceeds the bounded lower terms
and $E_*d$ is small.  Imaginary drift terms cost
$\varepsilon E|w_t|^2+C_\varepsilon Ed^2|w|^2$; the small mixed term
has the same bound with an additional small angular energy term.
Normalized coefficient derivatives cost only absorbable derivative
energies and bounded multiples of $|w|^2$.  For the boundary, the
condition is $w_t+dw=0$ at $T$.  The principal terminal radial and
drift contributions sum to
\[
 d\operatorname{Re}a\,|w(T)|^2
       +\tfrac12(1-2d\operatorname{Re}a)|w(T)|^2
 =\tfrac12|w(T)|^2.
\]
The $d_1$ term adds $\tfrac12\operatorname{Re}d_1(T)|w(T)|^2$,
which preserves this favorable sign for sufficiently small $E_*$.
Entrance Dirichlet data are first lifted.  Combining this coercive
estimate with Lemma~\ref{lin:strip} controls the complete bulk norm.
For the uniformly local version center at any strip and insert an
additional smooth weight comparable to
$\exp(-\epsilon|t-t_\nu|)$.  Its derivatives use less than the strict
rate buffer.  The squared source integrals are bounded by
$\sum_k e^{-c\epsilon|k-\nu|}\|f/\rho\|_{L^2(I_k;H^h)}^2$.
The geometric sum is independent of $T$.  This proves
\eqref{lin:forward} without a factor depending on the number of strips.
Smooth bounded interpolations between two admissible fixed rates add
$-a(\log\rho)''$, which is handled in the same estimate.
Finite entrance-Dirichlet/outgoing-derivative problems have index zero;
the estimate therefore also proves their solvability.  Exhaustion and
local compactness give the infinite exterior inverse.  Uniqueness in
the weighted uniformly local domain follows by the same localized
coercive estimate: remote cutoff terms vanish in the additional
exponentially decaying localization weight.  In particular no fast
outward homogeneous mode is admitted in that domain.

\begin{lemma}[Radial coefficient estimates]\label{gauge:radial}
In the cylindrical radius variable, the normalized stationary operator
is analytic from a sufficiently small strip graph neighborhood, with
the norm \eqref{lin:strip-norm} and traces \eqref{lin:first-traces},
to $L^2_tH^h_\theta$.  Its coefficient variables are
\[
 R,\quad R_\theta,\quad \mu R_t,
\]
and it is affine in the bulk second-derivative variables
\[
 \mu^2(R_{tt}-R_t),\quad \mu R_{t\theta},\quad
 R_{\theta\theta}.
\]
The graph neighborhood keeps the radius and metric denominators away
from zero and retains the strict principal accretivity margins.
If $v=R-\sqrt2$ is small in the unweighted uniformly local strip norm,
the mean-value factorization
\[
 \mathcal E_{\mathrm{rad}}(v)
 =(\mathcal L_{\mathrm{cyl}}+Q_v)v,\qquad
 Q_v=\int_0^1\!\bigl(D\mathcal E_{\mathrm{rad}}(sv)-D\mathcal E_{\mathrm{rad}}(0)\bigr)\,ds,
\]
of minus twice the radial graph velocity obeys
\[
 \|Q_v w\|_{Y_{\rho,h}}
 \le C_{\rho,h}\|v\|_{X_{1,h}}\|w\|_{X_{\rho,h}}
\]
for every fixed weight $\rho$ with bounded logarithmic derivatives.
The same coefficient rule applies after fixed real translations or
angular commutations when the corresponding graph norms are controlled.
\end{lemma}

\begin{proof}
Put $R_z=\mu R_t$ and
$D=R^2(1+R_z^2)+R_\theta^2$.  Direct inversion of
$g=\operatorname{diag}(1,R^2)+dR\otimes dR$ gives
\[
 g^{zz}=\frac{R^2+R_\theta^2}{D},\qquad
 g^{z\theta}=-\frac{R_zR_\theta}{D},\qquad
 g^{\theta\theta}=\frac{1+R_z^2}{D}.
\]
These are analytic functions of the displayed first-derivative
variables while $D$ is a unit.  In the curvature numerator above,
$R_{zz}=\mu^2(R_{tt}-R_t)$ and
$R_{z\theta}=\mu R_{t\theta}$; the remaining terms have only the
same first variables and $R_{\theta\theta}$.  The support term
$(R-R_t)/2$ is affine.  Thus every highest second derivative remains
in its $L^2_tH^h_\theta$ slot, while its coefficient is in
$L^\infty_tH^h_\theta$ by \eqref{lin:first-traces}.  The product estimate
\[
 \|AB\|_{L^2_tH^h_\theta}
 \le C_h\|A\|_{L^\infty_tH^h_\theta}
          \|B\|_{L^2_tH^h_\theta}
\]
and analytic inversion of the unit denominator prove the map assertion.
In its derivative, a coefficient variation is estimated in the first
trace norm and multiplies a background second derivative in the bulk
norm.  This proves the stated small mean-value bound, strip by strip;
division by $\rho$ gives the weighted version.  Fixed translations
preserve the factors $\mu^2,\mu$, and angular commutations use the
same algebra and module estimate.
\end{proof}

\begin{lemma}[Separated cylindrical rates]\label{lin:rates}
For \eqref{lin:cylinder-operator}, and for perturbations satisfying
the normalized decay condition below, a weight $e^{\alpha t}$
with $\alpha\notin\{2-j^2:j\in\mathbb Z\}$ gives separated end
estimates.  Slow data of rate below $\alpha$ are prescribed at the
entrance; slow data of rate above $\alpha$ are prescribed at infinity
(or at the appropriate finite exit).  Finite intervals have logarithmic
length at least a fixed $\ell_*>0$, with constants allowed to depend
on $\ell_*$.  Crossing a rate changes the
solution space by exactly the corresponding angular modes.  Estimates
hold with every fixed number of log-coordinate derivatives when the
source has the corresponding regularity, with an arbitrarily small
loss in a stated polynomial power.

Here the second-order differential perturbation
$Q=\mathcal L-\mathcal L_{\mathrm{cyl}}$ retains
the strict principal accretivity of Lemma~\ref{lin:strip}, with
$E\asymp e^{-2t}$, and, for the chosen $0<\eta<1$, has the following gain on every fixed weight
used in the assertion:
\begin{equation}\label{lin:normalized-decay}
 \|Qv\|_{Y_{e^{(\beta-2+\eta)t},h}}
       \le C_{\beta,h,\eta}\|v\|_{X_{e^{\beta t},h}}.
\end{equation}
The analogous bound is required after the fixed translations and
commutations used to obtain derivative estimates.  Thus the changes
of the radial, mixed, and angular principal coefficients are measured
in their respective units $E$, $\sqrt E$, and $1$; lower coefficient
terms may instead be measured in the corresponding coefficient--module
operator norm.  Ordinary absolute decay of the radial coefficient is
not the condition in \eqref{lin:normalized-decay}.

Consequently, on each cylindrical end and for every small $\eta>0$,
\begin{equation}\label{lin:end-expansion}
 R-\sqrt2=O(z^{-2+\eta}),\qquad
 J=c z^2+z(\ell\cdot e_r)+O(z^\eta)
\end{equation}
for a Gaussian Jacobi solution with compact source.  If $c=0$,
subtracting the exact infinitesimal rotation with tilt $\ell$ leaves
$O(z^{-2+\eta})$.  End solutions with arbitrary prescribed $c,\ell$
exist, and their inward cutoffs have compact Jacobi source.
\end{lemma}

\begin{proof}
High Fourier numbers satisfy the coercive estimate after any fixed
shift $\alpha$, so only finitely many angular modes need separate
treatment.  In such a mode the equation is
\[
 v_t-2\mu^2(v_{tt}-v_t)-\sigma_jv=f.
\]
Write $a=2\mu^2$, $q=v_t$, $k=(1+a)/a$, and
\[
 H(t,s)=\exp\!\left(-\int_t^s k(\zeta)\,d\zeta\right),\qquad
 (\mathsf A_T w)(t)=\int_t^T a(s)^{-1}H(t,s)w(s)\,ds.
\]
Solving the first-order equation for $q$ backward gives the exact formula
\[
 q(t)=H(t,T)q(T)+\mathsf A_T(\sigma_jv+f)(t).
\]
The kernel has bounded mass, width $O(\mu(t)^2)$, and an adjacent-strip
tail bounded by $C\exp(-c|s-t|/\mu(t)^2)$.  Its finite mass is
\begin{align*}
 \mathsf A_T1
 &=1-H(t,T)-\int_t^T\frac{a(s)}{1+a(s)}k(s)H(t,s)\,ds\\
 &=1-H(t,T)+O(a(t)).
\end{align*}
In particular the infinite average has mass $1+O(a(t))$, whereas the
finite one has a terminal layer.  On its last strip
$\|H(\cdot,T)\|_{L^2}\le C\sqrt{a(T)}=C\mu(T)$.
Translation across the kernel and the fundamental theorem of calculus
therefore give, for $\rho=e^{\alpha t}$,
\begin{equation}\label{lin:fast-average}
 \begin{split}
 \|(\mathsf A_T-I)v\|_{Y_{\rho,h}}
 \le{}&C\sup_{t\ge t_0}\mu(t)\,\|v_t\|_{Y_{\rho,h}}
       +C\sup_{t\ge t_0}\mu(t)^2\,\|v\|_{Y_{\rho,h}}\\
 &+C\mu(T)\rho(T)^{-1}\|v(T)\|_{H^h}.
 \end{split}
\end{equation}
One may equivalently retain $-H(t,T)v(T)$ as part of the terminal lift;
replacing $v(t)$ by $v(T)$ in that layer costs the first term above.
For the finitely many modes under consideration, the one-dimensional
$H^1$ trace on a terminal interval of length $\min(1,\ell_*)$ bounds
the last term by
$C\sup\mu(\|v\|_{Y_{\rho,h}}+\|v_t\|_{Y_{\rho,h}})$.
At an infinite exit the layer is absent.  The estimate thus retains a
small factor at every permitted finite exit, with no additional angular
trace requirement.
Keep this average on $f$.  The remaining first-order operator
$\partial_t-\sigma_j$ is inverted by integration from the entrance
if $\sigma_j<\alpha$, and by integration from the right if
$\sigma_j>\alpha$.  Its weighted kernel has mass at most
$|\alpha-\sigma_j|^{-1}$.  Taking $t_0$ large absorbs the displayed
averaging error.  The equation then controls $v_t$ and $\mu^2v_{tt}$.
This constructs the separated inverse, including finite exits.
Variation of constants shows that crossing $\sigma_j$ adds one solution
for each vector of its angular eigenspace, with that vector times
$e^{\sigma_jt}$ as leading term.  Condition
\eqref{lin:normalized-decay} makes the perturbation norm on a far tail
arbitrarily small at each fixed weight; the Neumann series and the same
variation-of-constants formulas then give the asserted perturbation
and crossing statements.

For the first assertion of \eqref{lin:end-expansion}, we first need
small-coefficient invertibility, without a coefficient decay hypothesis.
Put $v=R-\sqrt2$, $a=2\mu^2$, and let $K$ denote the exact curvature
numerator preceding \eqref{lin:cylinder-operator}.  Stationarity gives
\[
 q:=v_t=zR_z=R+2K\longrightarrow0,\qquad
 a v_{tt}=2(R_{zz}+R_z/z)\longrightarrow0,\qquad
 \sqrt a\,v_{t\theta}=\sqrt2R_{z\theta}\longrightarrow0.
\]
Indeed spatial convergence gives $K\to-1/\sqrt2$; angular
differentiation gives the same assertions at every fixed angular order.
Together with $v,v_{\theta\theta}\to0$, these identities prove that
the actual graph has small tail norm in $X_0:=X_{1,h}$, with all five
slots of \eqref{lin:strip-norm}.  No unweighted bound for $v_{tt}$
has been used.  The homogeneous term of the model fast derivative
inverse is proportional to $\exp(t+e^{2t}/4)$ and is excluded by
$q=zR_z=o(e^t)$.

Fix $0<\eta<1$, set $\alpha=-2+\eta$, and write
$X_\beta=X_{e^{\beta(t-t_0)},h}$ and similarly $Y_\beta$.
For both $\beta=0$ and $\beta=\alpha$, the model construction above
has the same entrance projection $S=P_{\{|j|\ge2\}}$ and gives
\[
 \|u\|_{X_\beta}\le C_{\eta,h}
 \bigl(\|\mathcal L_{\mathrm{cyl}}u\|_{Y_\beta}
       +\|Su(t_0)\|_{\mathcal T^D_{a(t_0),h}}\bigr).
\]
The other slow channels use their backward integrals at infinity;
every fast outward channel is excluded.  For the finitely many low
modes the integral kernel is bounded by the inverse spectral gap.
For $|j|\ge2$, conjugation by $e^{\beta(t-t_0)}$ gives interior energy
\[
 \int\!\left[a|w_t|^2+
 \{j^2-2+\beta-a(1+\beta+\beta^2)\}|w|^2\right].
\]
Here $a'=-2a$ and the leading mass is at least $\eta$.  Homogeneous
outgoing derivative data give terminal coefficient $(1-a(T))/2>0$.
The strip estimate supplies the remaining slots, and localization as
in \eqref{lin:forward} gives a uniformly local constant independent
of the terminal section.  Finite problems and exhaustion therefore
construct this inverse on each space, with the natural entrance trace
\eqref{lin:dirichlet-trace}.

Lemma~\ref{gauge:radial} supplies the nonlinear coefficient estimate
in these same strip norms, with the first-derivative traces and the
strict principal margins just established.

Factor the actual equation as $(\mathcal L_{\mathrm{cyl}}+Q)v=0$
by integrating its derivative along $sv$, $0\le s\le1$.
The radial coefficient variables $R,R_\theta,\mu R_t$, the first
traces \eqref{lin:first-traces}, and the coefficient--module rule
of Lemma~\ref{gauge:radial} give
$\|Q\|_{X_\beta\to Y_\beta}\le\delta(t_0)\to0$ for these two
weights.  Background second derivatives stay in their bulk slots;
no log derivative of a coefficient is needed here.
If $H_\beta$ lifts the stable entrance value and $K_\beta$ is the
zero-entrance model inverse, solve
\[
 u_\beta=H_\beta Sv(t_0)-K_\beta Q u_\beta
\]
by a Neumann series, choosing $C_{\eta,h}\delta(t_0)<1/2$.
The constructed $u_\alpha$ lies in $X_0$, as does the actual $v$.
Uniqueness in $X_0$ with the same stable data gives
$v=u_0=u_\alpha$.  This proves the domain upgrade and the first
bound in \eqref{lin:end-expansion}; it does not assume that bound
when estimating $Qv$.

To obtain log regularity, use the nonlinear equation itself.  Replace
its explicit $\mu(t)$ by $e^{-\tau}\mu(t)$ on the fixed tail and
prescribe stable entrance data $Sv(t_0+\tau)$.
The coefficient--module formulas make this a smooth family of maps
$X_\alpha\to Y_\alpha$, whose derivative in the unknown is a
small perturbation of the same invertible model.  The implicit function theorem
constructs its smooth local solution family.  The shifted actual graph
is small in $X_0$ and has these data.  The same mean-value estimate
gives uniqueness of small nonlinear solutions in $X_0$, identifying
it with that family without assuming translation differentiability
in $X_\alpha$.  Differentiating the explicit coefficients preserves
the factors $\mu^2$ and $\mu$ on second derivatives.  Iteration,
at a sufficiently large fixed angular order, gives
$\partial_t^k v\in X_\alpha$ for each fixed $k$; a small extra loss
may be allowed in $\eta$.  Consequently the normalized coefficient
differences and their required log derivatives decay like
$O(e^{-(2-\eta)t})$.  Only at this point do we use the
decaying-coefficient assertion for the sharper Jacobi asymptotics.

We also explain why the Gaussian condition permits use of polynomial
weights.  Local elliptic estimates on spatial balls of radius
$(1+|X|)^{-1}$ and the Gaussian $L^2$ bound give, for a homogeneous
exterior solution, a pointwise bound by a polynomial times
$\exp(|X|^2/8+o(|X|^2))$.  Choose $\alpha_0$ strictly between $1/8$
and $1/4$.  Direct differentiation gives
\[
 L e^{\alpha_0|X|^2}
  =\bigl((4\alpha_0^2-\alpha_0)|X|^2+O(1)\bigr)
                  e^{\alpha_0|X|^2}.
\]
Likewise $L(1+|X|^2)^{M/2}$ is negative on a far end for sufficiently
large fixed $M$, and dominates a given polynomial source.  Compare
the real and imaginary parts with a fixed multiple of this power plus
$\varepsilon e^{\alpha_0|X|^2}$.  The latter dominates the remote
boundary; after exhausting the end, let $\varepsilon\downarrow0$.
The comparison principle is justified by division by the positive
supersolution, which makes the zeroth-order coefficient nonpositive.
Thus Gaussian solutions of polynomial-source equations have polynomial
growth.  For a nonsmooth source in an exterior range space, first
subtract a cutoff of the exterior particular solution; the same
argument applies to the homogeneous remainder.

Starting with that polynomial bound, descend through weights above $2$,
then through $2$ and $1$.  The resulting coefficients are $c$ and
$\ell$ in \eqref{lin:end-expansion}.  Coefficient errors multiplying
$cz^2$ have size $O(z^\eta)$, explaining the remainder stated there.
When $c=0$, subtract an exact rotation field, whose leading term is
$z(\ell\cdot e_r)$.  No remaining slow rate lies between $0$ and
$-2+\eta$, giving the sharper remainder.  Conversely prescribe each
slow coefficient in the variation-of-constants construction on the
tail and multiply the resulting exact end solution by an inward
cutoff.  It has polynomial growth and compact source.

Finally log regularity does not require holomorphic spatial charts.
Pull the operator back by a small real translation of $t$, cut off
near the fixed entrance.  Its coefficients are differentiable in the
translation parameter in the same mapping norms.  Invert the
differentiated equation and use uniqueness to identify its derivative
with $\partial_t v$.  Iteration, increasing the fixed angular order if
needed, proves the asserted fixed-order log estimates.
\end{proof}

\subsection{Infinite conical ends}

On a conical end use a link graph on the unit sphere and $z=e^t$.
Its scalar variable $v$ is a link displacement; a bounded $v$ gives an
ambient normal displacement of size $O(z)$.  Division by the normal
velocity coefficient of link change gives
\[
 v_t-\mu^2\bigl(A^{ij}(t,\theta,v,Dv)D_{ij}v
                         +B(t,\theta,v,Dv)\bigr),\qquad \mu=e^{-t}.
\]
Thus all lower terms other than $v_t$ carry $\mu^2$.  The coefficients
of the real reference equation and their required log and angular
derivatives are bounded, and the diffusion matrix is uniformly
positive definite.  Smooth conical asymptotics give these bounds.

\begin{lemma}[Conical energy, traces, and exhaustion]\label{lin:conical}
Let
\begin{equation}\label{lin:conical-operator}
 \mathcal Cv=v_t-\mu^2
       (a v_{tt}+2b v_{t\theta}+c v_{\theta\theta}
                               +d v_t+qv_\theta+ev).
\end{equation}
Assume the real symmetric matrix
$\left(\begin{smallmatrix}a&b\\b&c\end{smallmatrix}\right)$ is at least
$c_*I$, and the coefficients and their required derivatives are
uniformly bounded.  There is $t_*$ depending only on these bounds such
that, for $t_0\ge t_*$, the entrance Dirichlet problem on
$[t_0,\infty)\times S^1$ is an isomorphism in the spaces
\begin{align}
 \|v\|_{X^{\mathrm{con}}_h}
 ={}&\sup_{t\ge t_0}
       \bigl(\|v(t)\|_{H^{h+1}}+\|v_t(t)\|_{H^h}\bigr)
       +\|\mu^{-1}v_t\|_{L^2H^h}
       +\|\mu D^2v\|_{L^2H^h},\label{lin:conical-norm}\\
 \|f\|_{Y^{\mathrm{con}}_h}
 ={}&\|\mu^{-1}f\|_{L^2H^h}.\label{lin:conical-source}
\end{align}
More explicitly,
\begin{equation}\label{lin:conical-estimate}
 \|v\|_{X^{\mathrm{con}}_h}
 \le C_h\bigl(\|\mathcal Cv\|_{Y^{\mathrm{con}}_h}
                  +\|v(t_0)\|_{H^{h+3/2}}\bigr).
\end{equation}
The same bound holds on $[t_0,T]$, $T\ge t_0+2$, with $v_t(T)=0$,
uniformly in $T$.
For the entrance term the constant may depend on the fixed $t_0$.
Small complex perturbations in the full domain-to-range operator norm
preserve these assertions.
\end{lemma}

\begin{proof}
First take real data, zero entrance value, and a finite interval with
zero terminal time derivative.  Put
\[
 U=\|\mu^{-1}v_t\|_2,\quad
 V^2=\|\mu v_{t\theta}\|_2^2+\|\mu v_{\theta\theta}\|_2^2,
 \quad F_0=\|f/\mu\|_2,\quad \mu_0=e^{-t_0},
 \quad \ell=dv_t+qv_\theta+ev.
\]
Multiplication of $\mathcal Cv=f$ by $-v_{\theta\theta}$, followed by
one integration in $t$ and one in $\theta$ for the $a$ term, gives
\begin{align}
 &\tfrac12\|v_\theta(T)\|_2^2
 +\int\mu^2
       (a v_{t\theta}^2+2b v_{t\theta}v_{\theta\theta}
                                      +c v_{\theta\theta}^2)
 \nonumber\\
 &\quad=-\int f v_{\theta\theta}
       -\int\mu^2a_\theta v_t v_{t\theta}
       +\int(\mu^2a)_t v_t v_{\theta\theta}
       -\int\mu^2\ell v_{\theta\theta}.\label{lin:conical-first}
\end{align}
The boundary term suppressed in this identity is
$-[\int\mu^2a v_t v_{\theta\theta}]_{t_0}^T$.  It vanishes because
$v_{\theta\theta}(t_0)=0$ and $v_t(T)=0$.

The second multiplier is $\mu^{-2}v_t$.  Integrating the $a$ term in
time and the mixed term in angle gives
\begin{align}
 U^2+\tfrac12\int_{S^1}a(t_0)v_t(t_0)^2
 ={}&\int(f/\mu)(v_t/\mu)+\int c v_{\theta\theta}v_t
       +\int\ell v_t\nonumber\\
 &-\int(\tfrac12a_t+b_\theta)v_t^2.
 \label{lin:conical-second}
\end{align}
In particular the entrance term is favorable and is generally nonzero.
The cross term $\int c v_{\theta\theta}v_t$ is bounded directly by
$CUV$.  Integrating it in time would unnecessarily introduce an
unweighted $c_t v_\theta^2$ term; no integrability of $c_t$ is assumed.

Let $H_0=\|\mu(qv_\theta+ev)\|_2$.  The two identities imply
\[
 c_*V^2\le F_0V+C\mu_0^2UV+H_0V,
 \qquad
 U^2\le F_0U+CUV+C\mu_0^2U^2+H_0U.
\]
Product integration, using the zero entrance value, gives
\[
 \sup_t\|v_\theta(t)\|_2^2\le2UV,
 \qquad
 \sup_t\|v(t)\|_2\le(\int_{t_0}^\infty\mu^2)^{1/2}U
                         =\mu_0 U/\sqrt2.
\]
Consequently
$H_0\le C\mu_0\sqrt{UV}+C\mu_0^2U$.
Add a sufficiently small fixed multiple of the second energy inequality
to the first, and then choose $\mu_0$ sufficiently small.  Young's
inequality absorbs the cross terms and yields
$U+V\le CF_0$ with a constant independent of $T$.
The equation recovers the missing derivative:
\[
 a\mu v_{tt}=\mu^{-1}v_t-2b\mu v_{t\theta}
             -c\mu v_{\theta\theta}-\mu\ell-f/\mu.
\]
Thus $W=\|\mu v_{tt}\|_2\le CF_0$.  Since $v_t(T)=0$,
product integration also gives
$\sup_t\|v_t(t)\|_2^2\le2UW$.
These are all the zeroth-angular-order norms in
\eqref{lin:conical-norm}.

Commute with $\partial_\theta^j$ and keep the same principal operator
on the left.  After division of the new source by $\mu$, the new
highest terms are
\[
 \mu(\partial_\theta^k A^{rs})
            D_{rs}\partial_\theta^{j-k}v,\qquad 1\le k\le j.
\]
Each involves a strictly lower angular commutation of the already
controlled second derivatives.  Induction in $j$ therefore proves
\eqref{lin:conical-estimate} at every fixed $h$, without a time
derivative of $f$.  Nonzero entrance data are removed by an extension
on the first fixed unit collar.  A Fourier extension
$\varphi_j e^{-(1+|j|)(t-t_0)}$, with a fixed cutoff, has commuted
$H^2$ norm controlled by $\|\varphi\|_{H^{h+3/2}}$; its forcing in
\eqref{lin:conical-source} has the same bound with a constant depending
on $t_0$.

For each finite $T$ this is a scalar elliptic problem with entrance
Dirichlet and transverse outgoing derivative boundary conditions.
Its index is zero, as follows by deforming the positive principal
matrix to the identity and the lower terms to zero on the finite
domain.  The estimate excludes a kernel, so the problem is surjective.
Exhaustion, weak compactness in the weighted bulk spaces, and local
elliptic compactness give an infinite-end solution with the same bound.

For uniqueness on the infinite end it is necessary to justify removal
of terminal terms for an arbitrary element of
\eqref{lin:conical-norm}.  Indeed
\[
 \int_{t_0}^\infty
   \left|\frac{d}{dt}\|v_t(t)\|_{H^h}^2\right|dt
 \le2\|\mu^{-1}v_t\|_{L^2H^h}\|\mu v_{tt}\|_{L^2H^h}<\infty.
\]
The weighted $L^2$ condition forces the limiting derivative trace to
be zero.  Also, after angular integration,
\[
 \int\mu^2a v_t v_{\theta\theta}
 =-\int\mu^2a v_{t\theta}v_\theta
       -\int\mu^2a_\theta v_tv_\theta.
\]
Choose terminal sections tending to infinity where
$\|\mu v_{t\theta}\|_{H^h}\to0$.  The other factors have bounded
traces, so this boundary term vanishes along those sections.  The
terminal angular energy in \eqref{lin:conical-first} remains
nonnegative.  The same two-energy argument now proves the estimate
and uniqueness on the infinite domain.  Finally all second derivatives,
including $\mu v_{tt}$, have been controlled; a small complex change
is therefore a small bounded perturbation of the full isomorphism and
is inverted by a Neumann series.
\end{proof}

\begin{lemma}[Gaussian solutions on a conical tail]\label{lin:gaussian-conical}
Let $J\in H^2_G(S)$ solve $LJ=f_c$ with $f_c\in L^2_G$ compactly
supported.  On a conical tail outside that support, its link displacement
belongs to $X^{\mathrm{con}}_h$ for every fixed $h$, and obeys the
entrance estimate \eqref{lin:conical-estimate}.  The same conclusion
holds after a conical source in $Y^{\mathrm{con}}_h$ has first been
removed by its zero-entrance particular solution.  These assertions
concern the real reference operator and its complex-valued solutions.
\end{lemma}

\begin{proof}
Fix a far entrance and write $z=e^t$.  The normal component of link
change is $z a_n(t,\theta)$, where $a_n$ is a real smooth unit with
bounded inverse and bounded fixed log and angular derivatives.  The
scalar normalization of the conical equation gives another such unit
$c_n$ for which, on the tail,
\[
 L(z a_n v)=z c_n\mathcal Cv.
\]
The functions $a_n,c_n$ include the normal-speed and drift factors;
they multiply the source as well as the unknown.  This identity follows
by differentiating the prescribed link graph equation, whose coefficient
of $v_t$ is a nonzero unit before its drift normalization.

Local elliptic regularity makes the trace
$v_0=J/(z a_n)$ at the entrance smooth.  Solve
$\mathcal Cv=0$ with this trace by Lemma~\ref{lin:conical}.  The normal
lift $\widetilde J=z a_n v$ is in the Gaussian $H^2$ domain on the tail.
Indeed smooth conical coordinates give $dA\asymp z^2dt\,d\theta$ and
\[
 |\nabla\widetilde J|\le C(|v|+|Dv|),\qquad
 |\nabla^2\widetilde J|
       \le Cz^{-1}(|D^2v|+|Dv|+|v|).
\]
The highest Gaussian integral is bounded by
\[
 \int e^{-z^2/4}|D^2v|^2\,dt\,d\theta
       \le C\int z^{-2}|D^2v|^2\,dt\,d\theta,
\]
and the other terms follow from the traces and weighted bulk terms in
\eqref{lin:conical-norm}.  A fixed collar extension at the entrance
preserves this conclusion.

The difference $J-\widetilde J$ is therefore a Gaussian exterior
$L$-solution with zero Dirichlet trace at the entrance.  Local elliptic
estimates give the pointwise bound used in the proof of
Lemma~\ref{lin:rates}, namely a polynomial times
$\exp(|X|^2/8+o(|X|^2))$.  For $1/8<\alpha_0<1/4$ the positive function
$e^{\alpha_0|X|^2}$ satisfies $Le^{\alpha_0|X|^2}<0$ on a sufficiently
far tail.  On each finite tail, divide the equation by this function
and apply the maximum principle to both signs of the real and imaginary
parts.  The entrance value is zero, and for each $\varepsilon>0$ the
remote value is bounded by $\varepsilon e^{\alpha_0|X|^2}$ once the
remote section is sufficiently far out.  Exhaustion and then
$\varepsilon\downarrow0$ give $J=\widetilde J$.  Thus the conical estimate
applies after, rather than before, identification of its domain.

For $f\in Y^{\mathrm{con}}_h$, its physical source $z c_n f$ belongs
to $L^2_G$, because
\[
 \int z^4|f|^2e^{-z^2/4}\,dt\,d\theta
       \le C\int z^2|f|^2\,dt\,d\theta.
\]
Solve the conical particular problem with zero entrance, cut its lift
off on a fixed collar, and subtract it.  The remaining source is
compact and the preceding argument applies.  The maximum principle
is used only for the real operator $L$; analytic perturbations in
complex parameters will be taken later on fixed real-coordinate
function spaces.
\end{proof}

\subsection{Compact norms and coefficient maps}

The infinite-end problems are joined to the Gaussian core on fixed
collars.  We keep the commuted compact norms explicit so that the
nonlinear map and the cutoff sources use the same domain and range.

Here is a fixed convention on the core and chart overlaps.  Choose a
finite family $\mathcal Z$ of smooth vector fields that spans the
tangent bundle on the core, contains the angular field on each end
collar, and reduces to angular fields sufficiently far out.  At least
one field is nonzero on every patch used in the transition.  On a
compact set $K$ put
\begin{equation}\label{lin:compact-norm}
 C_h^2(u;K)=\sum_{|\alpha|\le h}\|Z^\alpha u\|_{H^2(K)},
 \qquad C_h^0(f;K)=\sum_{|\alpha|\le h}\|Z^\alpha f\|_{L^2(K)}.
\end{equation}
The sums include words in the fields and the empty word.  Nested
compact collars are fixed once and for all.  Their elliptic estimates
control the entrance traces in Lemmas~\ref{lin:strip} and
\ref{lin:conical}; there is no artificial boundary at a change of
graph gauge.
For a smooth reference operator, $[L,Z]$ has order two with bounded
coefficients on these fixed sets.  Thus $[L,Z^\alpha]u$ is controlled
by the commuted $H^2$ norms of strictly smaller commutation order.
Interior estimates applied inductively to $Z^\alpha u$ justify the
claimed commuted estimates without loss of a derivative of the source.

\begin{lemma}[One graph chart across compact overlaps]\label{gauge:fiber-atlas}
Let $\vartheta$ denote the finite end-rotation parameters.  After fixing
real coordinate identifications on the compact overlaps, the core and
prescribed end graphs can be represented by one fiber map
\[
 \Gamma_\vartheta(y,s),\qquad \Gamma_0(y,0)=X(y),
\]
over the fixed real base atlas of $S$.  It is smooth in $y$ and analytic
in $(\vartheta,s)$.  It agrees with the fixed core graph map on the common
observation support, independently of $\vartheta$, and with the prescribed
rotated radius or link map sufficiently far on each end.  On the compact
overlaps, $(y,s)\mapsto\Gamma_\vartheta(y,s)$ is a local diffeomorphism
for small real $(\vartheta,s)$.  Its scalar derivative supplies the
linear graph-unit identifications on those overlaps.
\end{lemma}

\begin{proof}
Pull the core and end fiber maps once by the fixed real coordinate
identifications, so that their ambient values are functions of the same
base point $y$.  At $(\vartheta,s)=(0,0)$ all these values equal $X(y)$.
Orient their scalar variables so that the normal component of each scalar
derivative is positive.  On each fixed compact overlap interpolate the
ambient fiber maps pointwise with a smooth partition of unity in $y$.
At the reference, the base derivatives of the interpolated map are
$D_yX$, and its scalar derivative still has positive normal component.
The full derivative in $(y,s)$ is therefore invertible.  Compactness
preserves this transversality for small parameters.  Choose the partition
to select the core map on the observation support and the relevant end
map outside the overlap collars.  Because its
cutoffs depend only on the real base point, the interpolation is analytic
in the finite variables and invokes no parameter-dependent spatial
composition.  The far tail charts retain their already stated scaled
domains and units.
\end{proof}

\begin{lemma}[Conical and compact coefficient maps]\label{gauge:compact}
The normalized conical stationary operator is analytic from a
sufficiently small link-graph ball in $X^{\mathrm{con}}_h$ to
$Y^{\mathrm{con}}_h$.  On a sufficiently small neighborhood in fixed
compact graph charts, the stationary operator is analytic from the commuted $H^2$ norm
\eqref{lin:compact-norm} to its commuted $L^2$ range.  The exact fiber
chart of Lemma~\ref{gauge:fiber-atlas} gives the compact transitions.
The fixed family of commutations can span the core and
reduce to angular commutations sufficiently far out without loss of
derivatives.
\end{lemma}

\begin{proof}
After dividing the conical output by $\mu$, its equation is
\[
 \mu^{-1}v_t-\mu\bigl(A^{ij}(t,\theta,v,Dv)D_{ij}v
                                      +B(t,\theta,v,Dv)\bigr).
\]
The norm \eqref{lin:conical-norm} gives uniformly bounded
first-derivative algebra traces, which multiply $\mu D^2v$ in
$L^2_tH^h_\theta$.  The lower terms have the integrable factor
$\mu$.  Analytic composition on a small trace neighborhood and the
same algebra--module estimate as in Lemma~\ref{gauge:radial} therefore
prove the conical map assertion.
For example the trace of $v_t$ follows from product integration of
$\mu^{-1}v_t$ with $\mu v_{tt}$ on neighboring strips.

For completeness consider a compact flow box $(s,y)$ of one nonvanishing
commuting field, straightened to $\partial_s$.  Include
$\partial_s^j u$ in $H^2$ for $0\le j\le h$, and use overlapping
flow boxes for the finite field family.  In a differentiated coefficient
$A(u,Du)$, the sole potentially highest factor $\partial_s^h Du$ has
$L^\infty_yL^2_s$ control by its $y$ derivative and product integration.
The other factor $D^2u$ has $L^2_yL^\infty_s$ control by one additional
$s$ derivative.  Their product is in $L^2_{s,y}$, since
\[
 \|fg\|_{L^2_{s,y}}
 \le\|f\|_{L^\infty_yL^2_s}\,
       \|g\|_{L^2_yL^\infty_s}.
\]
When all commutations fall on $D^2u$, the coefficient has its ordinary
uniform first-derivative bound.  Intermediate distributions obey the
one-dimensional Sobolev product inequality in $s$, followed by this same
inequality in $y$.  Iterating the analytic coefficient expansion covers
any number of first-derivative factors.  Smooth changes of boxes and
cutoffs only create lower commutations.  This proves the required product
map on the compact transitions, including at a place where only one of
the retained fields is nonzero.
\end{proof}


\subsection{The analytic family preserving end types and its opening jets}

We now use the Gaussian and infinite-end inverses to separate changes that preserve end types from
quadratic cylindrical openings.  Analyticity is asserted for the former
parameters only.  Successive inversion constructs each coefficient in
the latter parameters, and Gaussian integration gives its action
coefficient; convergence of the opening series is not required.

\begin{proposition}[Normalized family and polynomial jets]\label{lin:family}
Let $\kappa$ be the number of cylindrical ends.  The observations can
be chosen so that their values are $q=\mathsf P(p)$ in analytic reduced
coordinates $p=(b,x)\in\mathbb R^{m-\kappa}\times\mathbb R^\kappa$,
with the following properties.
\begin{enumerate}
 \item At $x=0$ there is a real analytic family $S_b$ of normalized
 solutions of \eqref{lin:normalized-equation}.  It has the same end
 types as $S$, permits independent end tilts, and permits changes of
 its conical links within the normalized family.
 Outside the fixed observation support it is an exact shrinker.
 Its cylindrical remainder after rotating the end is
 $O(z^{-2+\eta})$ for every fixed $\eta>0$, on a parameter neighborhood
 that may depend on this choice.  On each conical end its link
 graph satisfies
 \begin{equation}\label{lin:family-conical-asymptotics}
 v_b(t,\theta)=v_{b,\infty}(\theta)+O(e^{-2t})
 \end{equation}
 in every fixed log and angular derivative, with $v_{b,\infty}$ analytic
 in $b$.  These bounds and their fixed parameter derivatives are uniform
 on a closed small parameter ball.
 \item There is a unique formal normalized solution in $x$,
 \[
 u(b,x)\sim\sum_{\alpha\in\mathbb N^\kappa}u_\alpha(b)x^\alpha,
 \qquad\lambda(b,x)\sim\sum_\alpha\lambda_\alpha(b)x^\alpha,
 \]
 based at $S_b$, whose coefficients in the prescribed graph charts
 have polynomial growth and are analytic in $b$.  For real $b$, their
 coefficientwise conversions to normal graph jets are smooth in $b$ and
 have uniform polynomial bounds at every fixed derivative order.
 After removal of the first-order tilt, the term homogeneous of degree
 one in $x$ on cylindrical end $j$ is
 \begin{equation}\label{lin:opening-coefficient}
  \frac{\sqrt2}{2}\,\beta_j(b,x)z^2+O(|x|z^\eta),
  \qquad \beta(b,x)=M(b)x,
 \end{equation}
 where $M(b)$ is an analytic invertible $\kappa\times\kappa$ matrix.
 In particular $\beta$ is exactly homogeneous linear in $x$ for fixed
 $b$.
 \item Integrating the Gaussian action coefficient by coefficient
 defines $\mathcal F\in\mathbb R\{b\}[[x]]$, with
 \begin{equation}\label{lin:action-multipliers}
  \nabla_p\mathcal F=(D_p\mathsf P)^{\mathsf T}\lambda.
 \end{equation}
 Every fixed truncation of the prescribed graph series and every fixed
 number of its parameter or spatial derivatives has uniform polynomial
 end bounds on a sufficiently small complex $b$ neighborhood.  The real
 normal truncations have the analogous smooth bounds for real parameters.
 This neighborhood may
 depend on the fixed truncation and derivative orders.
\end{enumerate}
When the limiting links and cylindrical directions are separated, as for
the tangent section in \Cref{geo:tangent}, the real family on a closed
sufficiently small ball has the following additional geometry.  Its smooth
rescaled conical annuli, the rotated cylindrical tails, and the compact
core have uniform bounded geometry.  There is a uniform positive normal
tubular radius; on a sufficiently far conical annulus of radius $r$ the tubular radius is
at least $cr$.  In the prescribed parametrizations
$|\partial_bX_b|\le C(1+|X|)$ and its physical first derivative is
bounded.  These statements concern the real surfaces; complex parameter
continuation is taken in the prescribed coordinates on a fixed real atlas.
\end{proposition}

\begin{proof}
\emph{The type-preserving inverse and the reference family.}
For each end, Lemma~\ref{lin:rates} constructs Jacobi solutions with
its opening and two tilt coefficients prescribed independently.
Extend them inward by cutoffs.  They lie in $H^2_G$ and have compact
Jacobi sources.  Apply Lemma~\ref{lin:observations} to these functions.
The map from the enlarged observation values to the $\kappa$ opening
coefficients of the homogeneous normalized linear solution is surjective.
Split the observation space into its kernel $V_b$ and a fixed
$\kappa$-dimensional complement $V_x$ on which this map is an isomorphism.

The type-preserving domain consists of the compact graph variables,
cutoff end rotations together with cylindrical remainders of weight
$-2+\eta$, and the conical link variables of
\eqref{lin:conical-norm}.  Its range has the corresponding cylindrical
source weight, conical source norm, and compact norms.  Impose only
the $V_b$ components of the observations, while retaining the full
multiplier vector.  We verify that the linearization is an isomorphism
in these spaces.  For a general exterior source, first solve a cylindrical
particular problem with zero slow coefficients at infinity and zero stable
entrance data in Lemma~\ref{lin:rates}, and a conical particular problem
with zero entrance value.  Cut their normal lifts off inside fixed far end
collars where the retained fields $\mathcal Z$ have reduced to angular
commutations.  Choose these collars outside the common compact observation
and multiplier-source support, so the cutoff lifts vanish on that support
and have zero observations.  Their Gaussian norms and those of their lifted sources are
controlled: this follows from their polynomial cylindrical bounds and
from the lift estimates in Lemma~\ref{lin:gaussian-conical}.

The remaining source is compact.  Solve the Gaussian normalized block
with the required $V_b$ observations and arbitrary $V_x$ observations,
then vary the latter uniquely until all opening coefficients vanish.
The finite-dimensional inverse used here is the opening map just
constructed.  Lemma~\ref{lin:rates} gives the resulting rotations and
decaying cylindrical remainders.  Enlarge the compact source set to include
the common multiplier-source support, and choose the later trace collars
outside this set.  Lemma~\ref{lin:gaussian-conical}, applied
at such an entrance, identifies the
conical link component with the solution in $X^{\mathrm{con}}_h$.
These estimates give a bound in the full type-preserving norm directly
for the stated range spaces.  On the core, the commuted compact elliptic
estimate bounds the $C_h^2$ norm by the Gaussian $H^2_G$ norm, the
commuted compact source norm, and the multiplier norm; the source columns
are fixed smooth compact functions.  The end inverses accept their full sources;
after the cutoffs the Gaussian source is compactly supported in $L^2_G$.
Near a farther fixed entrance its Gaussian solution is homogeneous and
therefore smooth, so the entrance trace estimates apply there even for
the original rough source.  A homogeneous element of this
domain lifts to $H^2_G$; if its $V_b$ observations and its openings vanish,
the $V_x$ observations also vanish.  The Gaussian block then kills the
element and its multipliers.  This proves the asserted isomorphism.

In prescribed graph coordinates the stationary equation is affine in
the highest second derivatives, with coefficients analytic in position
and first derivatives.  The denominator units are bounded away from
zero near the reference.  The trace and product rules above make it an
analytic map between precisely these domain and range spaces.  On a
conical end, for example, the divided range expression is
$\mu^{-1}v_t-\mu(A^{ij}D_{ij}v+B)$: bounded analytic coefficients of
the first-derivative traces multiply $\mu D^2v$ in $L^2$, and lower
terms have the integrable factor $\mu$.  The cylinder uses
$R,R_\theta,\mu R_t$ as coefficient variables and the second-derivative
slots in \eqref{lin:strip-norm}; compact transitions use
\eqref{lin:compact-norm}.  End rotations are prescribed ambient rotations
with a fixed cutoff on the core.  They are analytic in their finite
parameters without requiring analytic spatial charts.  The analytic
implicit function theorem gives $S_b$ and its multiplier vector
$\lambda_b$ on one parameter ball.  Write its full observation value
as $q_0(b)$ and set
\begin{equation}\label{lin:flattened-observations}
 \mathsf P(b,x)=q_0(b)+V_xx.
\end{equation}
This is a local analytic coordinate system for the observation values.
Write the preceding solution on the fixed atlas as
$X_b(y)=\Gamma_{\vartheta(b)}(y,u_b(y))$, and center its scalar chart by
\begin{equation}\label{lin:centered-fiber}
 \Psi_b(y,s)=\Gamma_{\vartheta(b)}(y,u_b(y)+s),\qquad
 \Psi_b(y,0)=X_b(y).
\end{equation}
This is the exact chart for subsequent prescribed graphs.  The low-space
family makes it analytic in the fixed-coordinate Banach spaces; the local
bootstrap below supplies every fixed spatial derivative.  On the
observation support the core fiber map remains fixed and each $\chi_i$
is evaluated at the same real $y$.  The variational source columns there
therefore acquire only the pointwise analytic density and normal-speed
units.  No spatial coordinate is composed with a $b$-dependent map.
For any smaller fixed cylindrical loss $\eta_0>0$, repeat this same
type-preserving inverse and implicit construction with weight
$-2+\eta_0$.  Its solution lies in the original weaker domain and has
the same $V_b$ observations.  Small uniqueness in that weaker domain
identifies it with $S_b$.  After shrinking the parameter ball for this
fixed loss and angular order, the same family is therefore analytic in
the stronger weight.  The same real-translation argument supplies its
fixed log derivatives there.  This justifies every fixed choice of the cylindrical
remainder loss used below.

\emph{Uniform regularity of the real family.}
The preceding construction supplies the small conical norm, but the
smooth geometry required below also uses the stationary equation.
Write the full link equation on a fixed tail as
\[
 \mathcal C^{\mathrm{nl}}(v)=v_t-\mu^2
 \{A^{ij}(t,\theta,v,Dv)D_{ij}v+B(t,\theta,v,Dv)\}=0.
\]
In the preceding construction choose one extra fixed entrance collar,
and analyze the tail starting at a later $t_0$.
On a fixed larger entrance collar, local parameter-dependent elliptic
bootstrap for this already constructed low-norm family gives every fixed
spatial Sobolev norm uniformly on a closed inner complex $b$-ball.
The coefficients and compact sources are smooth there and their ellipticity
is uniform.  The usual commuted local estimates prove the bounds successively;
testing against smooth functions and using the low-space holomorphy then
gives holomorphy in each of these locally bounded higher Sobolev norms.
In particular the entrance data used below are smooth in their real
translation parameters into $H^{h+3/2}$, uniformly on that ball.
This is a fixed-collar statement.

For small real translations $(\tau,\sigma)$ of the log and periodic angular
coordinates, prescribe on the fixed tail
\begin{align*}
 \mathcal C^{\mathrm{nl}}_{\tau,\sigma}(w)
 =w_t-e^{-2\tau}\mu(t)^2
 \{&A^{ij}(t+\tau,\theta+\sigma,w,Dw)D_{ij}w\\
   &+B(t+\tau,\theta+\sigma,w,Dw)\},\\
 w(t_0,\theta)&=v_b(t_0+\tau,\theta+\sigma).
\end{align*}
All explicit reference functions and the angular frame are shifted in this
formula.  The bounded reference derivatives and the conical coefficient--module
rule make it a $C^\infty$ family of maps from one fixed low
$X^{\mathrm{con}}_h$, $h\ge4$, to $Y^{\mathrm{con}}_h$, holomorphic in $b$.
Its derivative in $w$ is a uniformly small perturbation of the conical
entrance isomorphism.  The implicit function theorem therefore gives one
joint solution family, smooth in $(\tau,\sigma)$ and holomorphic in $b$.

The actual translated graph remains in the same small domain, even though
strong continuity of translation in the supremum part of that norm has
not been assumed.  On the enlarged tail,
\[
 \|\mu^{-1}v_{b,t}(\,\cdot+\tau)\|_2
 \le e^{|\tau|}\|\mu^{-1}v_{b,t}\|_2,\qquad
 \|\mu D^2v_b(\,\cdot+\tau)\|_2
 \le e^{|\tau|}\|\mu D^2v_b\|_2;
\]
the supremum bounds are unchanged up to restriction, and angular translation
is an isometry.  Small nonlinear uniqueness identifies the actual translate
with the implicit family.  Differentiating that one family gives
$\partial_t^k\partial_\theta^jv_b\in X^{\mathrm{con}}_h$ for every fixed
$j,k$.  Constants may depend on $j,k$, but these derivatives are taken on
one translation neighborhood and one closed inner $b$-ball.

The resulting first traces bound every fixed second log derivative
pointwise.  The equation now yields
$\|\partial_t v_b(t)\|_{H^h}\le C_h e^{-2t}$, and its differentiated
versions give the same factor for every fixed log, angular, and parameter
derivative of $\partial_t v_b$.  Integration to infinity proves
\eqref{lin:family-conical-asymptotics}, with
\[
 v_{b,\infty}=v_b(t_0)+\int_{t_0}^{\infty}v_{b,t}(t)\,dt.
\]
The integrable uniform estimates justify differentiation in $b$, so the
limiting link is analytic.  For the cylinder, the translation argument in
Lemma~\ref{lin:rates} also permits the angular translation $\sigma$:
its stable projection onto $|j|\ge2$ commutes with angular rotations.
Shift the explicit coefficients, angular frame, and stable entrance data
together, and use the same low-space uniqueness comparison with the rotated
family.  One joint translation family then gives all its fixed mixed
derivative bounds on a closed small parameter ball.

For real $b$, write a conical parametrization as $X_b=r\omega_b(t,\theta)$.
The preceding estimates give uniformly smooth rescaled annuli,
$|\mathrm{II}_b|\le C/r$, their fixed physical derivative bounds, and
$|X_b^\perp|\le C/r$ because $\omega_{b,t}=O(r^{-2})$.
The normalized metrics and their fixed derivatives vary by $O(|b|)$.
For the separated end data in the final assertion, the embedded limiting
links remain embedded and separated on a small closed ball.  Their tubular
radii on the unit sphere are uniformly positive.  Scaling gives a tube
of radius $cr$ on each far conical annulus; the annular asymptotics and
separated end directions prevent other ends from entering that tube.
The rotated cylindrical tails have a fixed positive tube, and compact
embeddedness gives one on the core and the fixed transition collars.
Taking the minimum proves the asserted global positive tubular radius.
A conical parameter variation is $r\partial_b\omega_b$ and a cylindrical
tilt also has size $O(r)$; their physical first derivatives are bounded.
This proves the final parameter-variation assertion.  In particular, on
a region of radius $R$, normal conversion between nearby real family
members costs at most $C(1+R)|b'-b|$ in height and slope while that quantity
is within the tube.

\emph{The augmented Gaussian inverse about the family.}
We establish this inverse before solving any opening coefficient.
Use the prescribed rotations and radial conical charts to pull all
differential expressions to the fixed real atlas of $S$.  Put $r=|X|$
there, and let
\[
 \mathcal D_G=\{w:\nabla^2w,\ (1+r)\nabla w,\ (1+r^2)w\in L^2(S)\}
\]
with its sum norm.  This is the half-density domain in
Lemma~\ref{lin:gaussian}.  For real $b$, let $F_b:S\to S_b$ be the
prescribed parametrization and $J_b$ its area Jacobian.  Pull back the
scalar component in the unit normal frame of $S_b$; if a prescribed graph
variable is used to compute that component, first multiply by its known
graph-speed factor.  Multiply the resulting physical normal scalar by
$J_b^{1/2}\exp(-(1+\delta)|F_b|^2/8)$.  This is the usual Gaussian
half-density identification for real $b,\delta$.

For complex $b$, take the analytic bilinear expressions for $F_b\cdot F_b$,
$J_b^{1/2}$, and the normalized normal frame on this same real atlas, and define
the conjugated differential expression by their chain rules.  The branches
are fixed from $b=0$.  This defines an analytic operator
$\widetilde{\mathcal A}_{b,\delta}:\mathcal D_G\oplus\mathbb C^m
\to L^2(S)\oplus\mathbb C^m$; it does not define a positive measure
on a complex surface or evaluate a spatial function at a complex point.
The uniform geometry just proved puts its coefficient difference from
$\widetilde{\mathcal A}_{0,0}$ on the three domain terms
\[
 o(1)|\nabla^2w|+
 o(1)(1+r)|\nabla w|+
 o(1)(1+r^2)|w|.
\]
For the complex coefficient bound, use the preceding joint translated
family as a holomorphic Banach-valued map at every fixed translated order.
The required first-trace maps are bounded; Cauchy's estimate on a closed
inner complex ball therefore gives $O(|b|)$ differences for all normalized
coefficient derivatives used here.  These are the analytic estimates
behind the real geometric interpretation.
Indeed normalized metric differences and their physical derivatives are
$O(|b|)$; differentiating the half-density contributes at most one factor
of $r$ in first-order coefficients and two in zeroth-order coefficients.
The measure-Jacobian derivatives have the bounded-geometry scales.
Changing $\delta$ contributes
$O(|\delta|)((1+r)|\nabla w|+(1+r^2)|w|)$.
The derivative of the base multiplier term
$\sum_i\lambda_{b,i}\chi_i(u)$ is retained and is a small compactly
supported operator, as are the observation and source-column changes.
Consequently
\begin{equation}\label{lin:gaussian-perturbation}
 \|(\widetilde{\mathcal A}_{b,\delta}
       -\widetilde{\mathcal A}_{0,0})(w,\lambda)\|
 \le \varepsilon(b,\delta)(\|w\|_{\mathcal D_G}+|\lambda|),
 \qquad \varepsilon(b,\delta)\longrightarrow0.
\end{equation}
The Gaussian block at $(0,0)$ is invertible.  A Neumann series therefore
gives a uniform inverse for small parameters, holomorphic in complex $b$,
on these fixed spaces.

For real $b,\delta$, its domain is also the derivative-defined Gaussian
$H^2$ domain on $S_b$ and controls the two moment terms.  To see this
directly, the compactly supported defect
$\Phi_b=\mathbf H_b+X^\perp/2$ gives the following divergence for the
weight $e^{-(1+\delta)|X|^2/4}$:
\[
 \operatorname{div}_{G,b,\delta}X^\top
 =2-\frac{1+\delta}{2}|X|^2+
       \frac{\delta}{2}|X^\perp|^2+\Phi_b\cdot X^\perp.
\]
The last two terms are uniformly bounded.  The moment integrations in
Lemma~\ref{lin:gaussian} therefore hold uniformly for this density.
Together with \eqref{lin:gaussian-perturbation} and its half-density
domain, they identify the domain and give its $|X|\nabla u$ and
$|X|^2u$ bounds.  This establishes the real shifted estimate recorded
after this proof, with the derivative of the nonzero base multiplier
included.

\emph{The mixed inverse for a fixed opening degree.}
Fix $h\ge4$ and a cylindrical rate $\gamma>2$.  Choose the tail entrances
far enough for this pair of parameters, absorbing the intervening finite
collars into the compact norms.  On the fixed real atlas
let $\mathcal X_{\gamma,h}$ consist of the compact norm
\eqref{lin:compact-norm}, the cylindrical norms
$X_{e^{\gamma(t-t_0)},h}$, and the original conical norms
$X^{\mathrm{con}}_h$, with the linear graph-unit identifications obtained
by differentiating \eqref{lin:centered-fiber} on fixed collars.
Let $\mathcal Y_{\gamma,h}$ have the
corresponding compact, cylindrical, and conical source norms.
The full augmented derivative about $S_b$, written in these units and
including the base compact multiplier derivative, satisfies
\begin{equation}\label{lin:mixed-polynomial-inverse}
 \mathcal A_b:
 \mathcal X_{\gamma,h}\oplus\mathbb C^m
       \longrightarrow\mathcal Y_{\gamma,h}\oplus\mathbb C^m
 \quad\hbox{is an isomorphism for small $b$, with analytic inverse.}
\end{equation}
The parameter neighborhood and its inverse bound may depend on
$\gamma,h$.  All spaces in this assertion are fixed real-coordinate
spaces; the conical part is the bounded-link domain, not a growing-link
space.

Here is the proof.  At $b=0$, solve a zero-entrance exterior particular
problem on each cylinder using Lemma~\ref{lin:rates} at $\gamma>2$,
where every slow channel is an entrance channel.  Solve the original
zero-entrance problem of Lemma~\ref{lin:conical} on each cone.  Cut their
normal lifts off on fixed far collars where $\mathcal Z$ contains only
the retained angular commutations, and denote the result by $u_{\rm ext}$.
Let $f_{\rm phys}$ be the prescribed source in physical normal units,
with the graph and row factors applied, and put
$r_c=f_{\rm phys}-Lu_{\rm ext}$.
The exterior estimates and the fixed cutoff product estimates give
\[
 \|u_{\rm ext}\|_{\mathcal X_{\gamma,h}}
 +\|u_{\rm ext}\|_{H^2_G}
 +\|r_c\|_{C_h^0(K)}
 +\|r_c\|_{L^2_G}
 +|Pu_{\rm ext}|
 \le C_{\gamma,h}\|f\|_{\mathcal Y_{\gamma,h}},
\]
where $K$ is one fixed compact set containing the remainder support and
the common compact observation and multiplier-source support.  The
Gaussian observation datum is the prescribed observation datum minus
$Pu_{\rm ext}$.  The Gaussian lift bounds on
cylinders follow by summing their polynomial strip bounds against the
Gaussian, and the conical lift and source bounds are those proved in
Lemma~\ref{lin:gaussian-conical}.  Solve this remaining compact source
and the remaining observations by the Gaussian augmented block.
Its solution $(J_G,\lambda_G)$ is bounded in $H^2_G\oplus\mathbb C^m$.
On fixed nested compacts containing $K$, the commuted elliptic estimate
applied to $LJ_G=r_c-P^*\lambda_G$ gives
\[
 C_h^2(J_G;K_{\mathrm{core}})
 \le C\bigl(\|J_G\|_{H^2_G}+\|r_c\|_{C_h^0(K)}+|\lambda_G|\bigr).
\]
The smooth compact source columns account for the last term, which the
Gaussian inverse already controls.
On nested fixed collars outside $K$, the commuted local elliptic estimates
bound the cylindrical $\mathcal T^D$ entrance traces and conical
$H^{h+3/2}$ entrance traces by
$C(\|J_G\|_{H^2_G}+\|r_c\|_{C_h^0(K)})$.
The polynomial Gaussian barrier and the separated estimate at $\gamma>2$
now bound the entire cylindrical component by those traces; equivalently
its expansion is \eqref{lin:end-expansion}.  On a cone,
Lemma~\ref{lin:gaussian-conical} identifies $J_G$ with the conical entrance
solution, whose full norm is bounded by its trace.  These inequalities
give the claimed real inverse bound for the full source spaces.  In
particular the cutoff source may be only compactly supported $L^2_G$;
homogeneity near the farther entrance is what gives the smooth traces.
Conversely an element of
$\mathcal X_{\gamma,h}$ has polynomial cylindrical local Sobolev bounds,
and its conical normal lift lies in $H^2_G$ by
Lemma~\ref{lin:gaussian-conical}'s lift calculation.  Its Gaussian norm
is finite, so the Gaussian block proves injectivity.

For complex $b$ close to zero, keep exactly the same link displacement
as domain variable on each conical tail.  The full drift normalization is
already part of $\mathcal C^{\mathrm{nl}}$, so its derivative at $v_b$
is directly a small analytic bounded perturbation from the fixed
$X^{\mathrm{con}}_h$ to $Y^{\mathrm{con}}_h$.  No global multiplication
of the conical domain by a graph-speed factor is made.  A further scalar
row factor, if used, acts only on $Y^{\mathrm{con}}_h$ and is an analytic
bounded $L^\infty_tH^h_\theta$ multiplier.  The factors $z a_n$ were used
for the real Gaussian lift above, and on fixed collars their conversions
are bounded in the compact norms.  On each cylinder the
prescribed rotation removes the tilt, and the small decaying remainder
gives the same small normalized coefficient--module perturbation on the
fixed weight $\gamma$.  The constant can depend on that weight.
The cylindrical row and graph-speed ratios used to identify the $b$ and
$0$ equations, and the scalar identifications on fixed collars, are
analytic bounded multiplication maps in their strip or compact modules;
all spatial coordinates remain the fixed real ones.  On the compact
support, the coefficients, base multiplier derivative, source columns,
and observation maps are analytic bounded maps in their fixed commuted
norms.  It follows that
\[
 b\longmapsto\mathcal A_b\in\mathcal B(
 \mathcal X_{\gamma,h}\oplus\mathbb C^m,
 \mathcal Y_{\gamma,h}\oplus\mathbb C^m)
\]
is holomorphic in the full operator norm, including both finite-dimensional
rows and columns, and $\|\mathcal A_b-\mathcal A_0\|\le C_{\gamma,h}|b|$
on a smaller ball.
A Neumann series proves \eqref{lin:mixed-polynomial-inverse}, including
its analytic complex parameter dependence.  This step supplies a
polynomial norm for the coefficients; Gaussian-space analyticity alone
would not supply that norm.

At degree one, choose $\gamma=2+\epsilon$ with $\epsilon>0$ small and
apply this inverse to the prescribed $V_x$ observations.  On a cylindrical
tail write its equation as
$\mathcal L_{\mathrm{cyl}}u_1=-Q_bu_1$, where the normalized gain
\eqref{lin:normalized-decay} for the rotated family is holomorphic in
$b$ in its full operator norm.  Choose its loss $\eta_0>0$ so that
$\epsilon+\eta_0<1$, using the corresponding stronger realization of
$S_b$ on its smaller parameter ball.  Then $Q_bu_1$ belongs holomorphically to
$Y_{e^{(\epsilon+\eta_0)(t-t_0)},h}$.  The fixed model
variation-of-constants splitting is a bounded linear map from this source
and the entrance trace of $u_1$ to the rate-$2$ coefficient, the two
rate-$1$ coefficients, and a remainder of rate $\epsilon+\eta_0$.
Every fast outward mode is excluded by the polynomial domain.
The trace map is bounded on $\mathcal X_{\gamma,h}$, so this splitting
proves analytic dependence of the opening and tilt coefficients without
assuming them in advance.  The two small losses can be chosen below any
specified remainder loss in \eqref{lin:opening-coefficient}.
At $b=0$ the resulting opening matrix is the isomorphism on $V_x$,
so it stays invertible.  With the normalization $c_j=\beta_j/\sqrt2$
it is $M(b)$ in \eqref{lin:opening-coefficient}.  Removing the
first-order tilt gives the displayed remainder.

Now suppose all opening coefficients of degree below $n$ have been
constructed in prescribed charts.  Their cylindrical source is controlled
directly by the weighted coefficient--module rule, before any pointwise
regularity bootstrap.  On a strip put
\[
 \mathcal J_1v=(v,v_\theta,\mu v_t),\qquad
 \mathcal J_2v=(\mu^2(v_{tt}-v_t),\mu v_{t\theta},v_{\theta\theta}).
\]
For $v\in X_{\rho,h}$ the first tuple has algebra trace norm at most
$C\rho\|v\|_{X_{\rho,h}}$, and the second has bulk norm at most the same
quantity.  Each formal coefficient of the radial equation is a finite sum
of products of first tuples and at most one second tuple, with bounded
analytic background coefficients.  The product rule therefore puts a
monomial with input weights $\rho_1,\ldots,\rho_j$ in the source weight
$\rho_1\cdots\rho_j$, with norm bounded by the product of the input norms.
The lower terms use one first-derivative factor in bulk and obey the same
bound.  Choose a fixed $\gamma>2$ above the finitely many resulting power
exponents, allowing a small extra power for logarithms.
On a conical end, keep every coefficient in the link variable.
The conical source at degree $\alpha_0$ is a finite sum of terms
\begin{equation}\label{lin:conical-jet-source}
 -\frac1{j!}D^j\mathcal C^{\mathrm{nl}}(v_b)
       [v_{\alpha_1},\ldots,v_{\alpha_j}],
 \qquad j\ge2,\quad
 \alpha_1+\cdots+\alpha_j=\alpha_0,\quad |\alpha_i|\ge1.
\end{equation}
The analytic map
$\mathcal C^{\mathrm{nl}}:X^{\mathrm{con}}_h\to Y^{\mathrm{con}}_h$ has bounded
multilinear derivatives on a small ball.  Every term in
\eqref{lin:conical-jet-source} therefore belongs to
$Y^{\mathrm{con}}_h$, with a bound by the product of the lower
$X^{\mathrm{con}}_h$ norms, and is analytic in $b$.
The compact source has the analogous finite Taylor bound.
Equation \eqref{lin:mixed-polynomial-inverse}, with the prescribed
coefficient of \eqref{lin:flattened-observations}, now gives the unique
coefficient of degree $n$.  In particular its conical link coefficient
again belongs to $X^{\mathrm{con}}_h$.  The chosen cylindrical weight
and the small complex $b$-ball may change with $n$.  This is precisely
the coefficientwise analyticity required by $\mathbb R\{b\}[[x]]$.

The following local argument closes the derivative part of this induction
at each degree.  Cylindrical bounds follow from the
commuted separated estimates, increasing the weight and angular order
by a fixed amount when needed.  On a conical unit log strip at height
$z$, the $X^{\mathrm{con}}_h$ norm initially bounds the local second
derivatives by a fixed multiple of $z$, and bounds the first traces.
After division by $\mu^2$, a coefficient equation has uniformly elliptic
principal matrix and a first-order drift of size $O(z^2)$; all fixed
derivatives of its other coefficients are bounded by the family
regularity already proved.  Commuted local elliptic estimates on nested
fixed log patches, with interpolation to absorb the first-order term,
therefore have constants bounded by a fixed power of $z$ at every fixed
order.  This follows inductively because each commutation differentiates
a coefficient a finite number of times, and the only growing coefficient
and its derivatives are $O(z^2)$.  The source at degree $n$ is a finite
product of lower-degree derivatives.  Starting from the preceding local
$H^2$ bound, induction on degree and on the local elliptic order gives
polynomial local Sobolev bounds for every fixed derivative.  Sobolev
embedding gives the corresponding pointwise bounds.  These bounds are
uniform on a closed smaller degree-dependent complex $b$-ball; fixed
coordinate differentiation and locally bounded holomorphy give the
same statement for parameter derivatives.

The conversion to normal jets is coefficientwise.  On a real scaled
annulus choose a local projection branch and use its normal chart to write the equality between the
prescribed graph parametrization and a normal graph with an unknown
real tangential displacement.  At $x=0$ this displacement is zero.
At a fixed $x$ degree, expand that equality at the identity using only
the finitely many spatial derivatives of the lower jets and of $S_b$
which occur at that degree.  The pointwise matrix formed by the tangent
vectors and the normal is invertible.  On a conical annulus, scale the
ambient and normal variables by its radius; the matrix and its inverse
are then uniformly bounded in the two link coordinates.  On a cylindrical
tail the axial and angular tangent scales are different in log coordinates:
the axial tangent has size comparable to $z$ and the angular tangent has
size comparable to one.
The prescribed cylinder charts therefore give an inverse with at most a
fixed polynomial factor in $z$; on compact charts it is uniformly bounded.
Solving this matrix equation gives the
tangential coefficient and the normal coefficient at that degree.
The local polynomial bounds just proved give polynomial bounds for
these coefficients and all fixed derivatives.  Differentiating the finite
real identities in $b$ gives the same fixed-order polynomial bounds for
their parameter derivatives.  For real parameters the
ordinary smooth implicit function theorem and Taylor's theorem on the
chosen projection branch identify these finite coefficients and the corresponding remainder
with the actual normal conversion.  More quantitatively, let
$v_M^{\rm pres}(b,x)$ be the finite prescribed graph polynomial through
degree $M$, let $v_M^{\rm nor}(b,x)$ be its coefficientwise normal
polynomial, and write $\mathfrak n_b$ for the actual normal-height
conversion for real parameters.  Wherever its scalar values lie in the
chart domain, the first polynomial represents the exact immersion
$y\mapsto\Psi_b(y,v_M^{\rm pres}(b,x)(y))$; its fixed-overlap
identifications agree exactly by \eqref{lin:centered-fiber}.
For the separated family, work on an inner
patch whose inverse projection stays in a controlled larger collar on the
chosen branch.  Suppose that throughout this larger collar, for every parameter
on the segment $\{tx:0\le t\le1\}$, the physical graph displacement
stays within a fixed fraction of the local tubular radius and its first
derivative is correspondingly small in the reference metric.  Then the
projection derivatives through any chosen finite order have at most
a fixed polynomial cost in $r=|X|$.  Taylor's theorem in the finite
variable $x$, applied to the smooth real projection equation, therefore
gives for every fixed spatial order $j$
\begin{equation}\label{lin:normal-conversion-remainder}
 \left|\nabla_{S_b}^j\bigl(\mathfrak n_b(v_M^{\rm pres}(b,x))
                   -v_M^{\rm nor}(b,x)\bigr)\right|
 \le C_{M,j}|x|^{M+1}(1+r)^{D_{M,j}}.
\end{equation}
The constants are uniform on a closed smaller real $b$-ball and on such
inner patches with the stated collar and tubular controls.  The complex parameter
assertions are made for the prescribed-coordinate jets, where they were
proved by the mixed inverse; the action is computed in those coordinates.

\emph{The formal action.}
Each fixed action coefficient is a polynomially bounded combination
of the finite graph jets and their first derivatives times the Gaussian.
On a sufficiently small degree-dependent complex $b$-ball, the prescribed
tilted and conical immersions satisfy
$\operatorname{Re}(F_b\cdot F_b)\ge c r^2-C$.
Gaussian domination permits coefficientwise integration and differentiation
in $b$.  In the formal first variation on exhausted domains, the boundary
terms are polynomially bounded times a Gaussian and tend to zero.
The plus normalization in \eqref{lin:normalized-equation} gives
$d\mathcal F=\lambda\cdot d\mathsf P$, proving
\eqref{lin:action-multipliers}.  The same bounds control every fixed
truncation.  The augmented coefficient equations are unique in the
mixed spaces.  To see uniqueness among all formal solutions with polynomial
growth, take the first degree at which two such solutions differ.  Their
lower sources and observation coefficients agree, so the difference solves
the homogeneous full augmented derivative about $S_b$.  At this degree
its prescribed graph difference converts linearly to a polynomially growing
physical normal scalar on the fixed real atlas.  Local elliptic estimates
for the homogeneous equation give polynomial bounds for its first two
derivatives as well.  After the analytic half-density conjugation it
belongs to $\mathcal D_G$, also for small complex $b$.  The full augmented
Gaussian inverse proved before the jets kills this difference and its
multiplier.  Induction proves uniqueness in the stated polynomial class.
For real parameters the finite normal-coordinate identities above identify
the same geometric jets; their action coefficients are those computed in
the prescribed charts.
\end{proof}

\begin{corollary}[Uniform shifted Gaussian comparison]\label{lin:shifted-gaussian}
Let $L_b^{\mathrm{norm}}$ denote the derivative of
\eqref{lin:normalized-equation} about $S_b$, including derivatives of
its compact multiplier terms.  For real $b$ sufficiently small and
real $|\delta|$ sufficiently small, put
$d\gamma_{b,\delta}=\exp(-(1+\delta)|X|^2/4)dA_{S_b}$.
There is a constant independent of these parameters such that
\begin{equation}\label{lin:shifted-estimate}
 \|u\|_{H^2_{G,b,\delta}}+|\lambda|
 \le C\left(
  \left\|L_b^{\mathrm{norm}}u+\sum_i\lambda_i\chi_{i,b}
        \right\|_{L^2_{G,b,\delta}}+|D P_bu|\right).
\end{equation}
The domain also controls $|X|\nabla u$ and $|X|^2u$ in this weight.
For complex $b$, the corresponding inverse is the analytic inverse
on the fixed pulled-back half-density spaces described in the proof
of Proposition~\ref{lin:family}.
\end{corollary}

\begin{proof}
The inverse was established before the opening-jet construction in
\eqref{lin:gaussian-perturbation}.  For real parameters its half-density
identification is unitary in the displayed positive measure, and the
moment calculation there identifies its domain with
$H^2_{G,b,\delta}$ and gives the two additional moment bounds.
Conjugating that inverse back proves \eqref{lin:shifted-estimate}.
All parameter perturbations are fixed before any later localization
radius is allowed to increase.
\end{proof}
